% Resumo: 150 a 250 palavras (regra do formato LNCS).
\begin{abstract}
Um veículo autónomo tem de prever o movimento dos outros utentes da via e
de usar essa previsão para planear o seu próprio movimento. Este artigo
revê o estado da arte nas duas frentes, a previsão de trajetórias e o
planeamento de movimento, e na forma como se ligam e se avaliam, com
atenção ao que conseguem os modelos de aprendizagem automática clássica,
sem redes neuronais. Apresenta ainda resultados preliminares de uma
arquitetura «prever para planear» sobre o Waymo Open Motion Dataset:
florestas aleatórias que aprendem, no referencial de cada veículo, a
correção a uma extrapolação de velocidade constante, e um planeador de
velocidade A* que trata as previsões como obstáculos móveis. Em 149 cenas
reais (13\,856 amostras, validação cruzada agrupada por cena), o erro
médio de posição desce de \un{4,97}{m} para \un{2,47}{m} a \un{3}{s} e de
\un{12,57}{m} para \un{6,66}{m} a \un{5}{s}, face à velocidade constante.
No planeamento, a versão atual do planeador ainda não converte esse ganho
numa redução mensurável das colisões; identificam-se as causas e o
trabalho necessário para o conseguir.

\keywords{\palavraschave}
\end{abstract}
